\documentclass[12pt, a4paper]{report}
\usepackage[utf8]{inputenc}
\usepackage[spanish]{babel}
\usepackage{graphicx}
\usepackage{hyperref}
\usepackage{geometry}
\geometry{a4paper, margin=2.5cm}

\title{\Huge\textbf{Título de tu TFG} \\ \vspace{0.5cm} \Large Despliegue y Análisis de una Red SDN}
\author{\textbf{Héctor}}
\date{\vspace{1cm} \today}

\begin{document}

% Portada
\maketitle

% Índice
\tableofcontents
\newpage

\chapter{Introducción}
Aquí puedes redactar la introducción de tu Trabajo de Fin de Grado. Explicarás la motivación, los objetivos (generales y específicos) y la estructura del documento.

\chapter{Estado del Arte y Tecnologías}
En este capítulo puedes hablar sobre los conceptos teóricos:
\begin{itemize}
    \item Redes Definidas por Software (SDN).
    \item El controlador Ryu.
    \item El emulador Mininet.
    \item Docker y contenedores.
\end{itemize}

\chapter{Arquitectura e Implementación}
Aquí describirás cómo has implementado el proyecto. Puedes incluir diagramas y explicar el funcionamiento del Firewall, el Dashboard y la API.

\chapter{Pruebas y Resultados}
Aquí mostrarás las pruebas de conectividad (Ping), los test de microsegmentación y las métricas representadas en tu Dashboard.

\chapter{Conclusiones y Trabajo Futuro}
Conclusiones extraídas tras la realización del proyecto y posibles mejoras futuras.

% Bibliografía básica (puedes usar BibTeX más adelante)
\begin{thebibliography}{9}
\bibitem{ryu} 
Documentación oficial de Ryu, \textit{https://ryu-sdn.org/}.
\bibitem{mininet} 
Documentación de Mininet, \textit{http://mininet.org/}.
\end{thebibliography}

\end{document}